\documentclass[12pt]{article}
\usepackage[margin=0.75in,top=0.8in,bottom=0.65in]{geometry}
\usepackage{lmodern}
\usepackage{microtype}
\usepackage{fancyhdr}
\usepackage{titlesec}
\usepackage{enumitem}
\usepackage{lastpage}
\usepackage[hidelinks]{hyperref}

\setlength{\parindent}{0pt}
\setlength{\parskip}{0.38em}
\setlength{\headheight}{26pt}
\titleformat{\section}{\large\bfseries\scshape}{}{0pt}{}
\titleformat{\subsection}{\normalsize\bfseries}{}{0pt}{}
\pagestyle{fancy}
\fancyhf{}
\fancyhead[C]{\footnotesize Lit Example Reader \quad Assignment ID: U1L2LE \quad Page \thepage\ of \pageref{LastPage}}
\renewcommand{\headrulewidth}{0.5pt}

\newenvironment{playpassage}{\begin{quote}\small\setlength{\parskip}{0.25em}}{\end{quote}}
\newcommand{\speaker}[1]{\textbf{#1.}}

\begin{document}

\begin{center}
{\LARGE\bfseries Week 2 Lit Example Reader}\par\vspace{0.05em}
{\large\scshape Macbeth: Shifting Words, Shifting Reasons}\par\vspace{0.15em}
\rule{0.78\textwidth}{0.5pt}
\end{center}

\textbf{Purpose:} This reader gives Macbeth passages for applying Week 2's logic habit: define your terms, watch for ambiguity, and do not let a word change meaning inside an argument. Read the passages as first-year college students: slowly, precisely, and with special attention to repeated words.

\textbf{What to watch for:} A word can be powerful in poetry and still dangerous in reasoning. Shakespeare often lets a word carry more than one sense at once. Logic does not flatten the poetry; it asks whether an inference depends on one sense of a word quietly becoming another.

\section*{Before the Passages: The Week 2 Logic Tool}

Whately's second habit is to keep terms steady. When an argument uses the same word in two senses, the reason may appear stronger than it is. For each passage, ask:

\begin{enumerate}[leftmargin=1.5em,itemsep=0.2em]
\item \textbf{Term:} Which important word or phrase is doing logical work?
\item \textbf{First sense:} What does it mean in one place?
\item \textbf{Second sense:} Does it mean something different in another place?
\item \textbf{Inference check:} Does the argument still work when the two senses are kept separate?
\end{enumerate}

Do not hunt for a fallacy label first. Begin with the words on the page. If the term stays steady, say so. If it shifts, explain exactly how.

\section*{Passage 1: Act 1, Scenes 1 and 3 --- Fair and Foul}

\textbf{Context:} The witches open the play by joining opposites. Later, Macbeth echoes their language before he even hears their prophecy. The repeated language prepares the reader for a world where attractive things may be morally dangerous and ugly things may appear successful.

\begin{playpassage}
\speaker{Witches} Fair is foul, and foul is fair:\par
Hover through the fog and filthy air.

\speaker{Macbeth} So foul and fair a day I have not seen.
\end{playpassage}

\subsection*{Reader Notes}

The words \textit{fair} and \textit{foul} can name more than one kind of judgment. \textit{Fair} may mean beautiful, favorable, honorable, or seemingly good. \textit{Foul} may mean ugly, filthy, morally corrupt, or unfavorable. In poetry, Shakespeare can press these meanings together. In logic, however, you must ask whether a claim about one sense proves a claim about another sense. A favorable outcome is not automatically an honorable one. An ugly battlefield victory is not automatically a morally foul action.

\textbf{Reading task:} Mark every use of \textit{fair} and \textit{foul}. Beside each one, write the most likely sense in that line.

\newpage

\section*{Passage 2: Act 1, Scene 7 --- What Counts as a Man?}

\textbf{Context:} Macbeth has decided not to murder Duncan. Lady Macbeth attacks his resolve by turning the word \textit{man} into a test. Her language matters because she is not merely describing Macbeth; she is pressuring him to accept a definition.

\begin{playpassage}
\speaker{Macbeth} We will proceed no further in this business:\par
He hath honour'd me of late; and I have bought\par
Golden opinions from all sorts of people,\par
Which would be worn now in their newest gloss,\par
Not cast aside so soon.

\speaker{Lady Macbeth} Was the hope drunk\par
Wherein you dress'd yourself? hath it slept since?\par
And wakes it now, to look so green and pale\par
At what it did so freely? From this time\par
Such I account thy love. Art thou afeard\par
To be the same in thine own act and valour\par
As thou art in desire? Wouldst thou have that\par
Which thou esteem'st the ornament of life,\par
And live a coward in thine own esteem,\par
Letting ``I dare not'' wait upon ``I would,''\par
Like the poor cat i' the adage?

\speaker{Macbeth} Prithee, peace:\par
I dare do all that may become a man;\par
Who dares do more is none.

\speaker{Lady Macbeth} What beast was't, then,\par
That made you break this enterprise to me?\par
When you durst do it, then you were a man;\par
And, to be more than what you were, you would\par
Be so much more the man.
\end{playpassage}

\subsection*{Reader Notes}

Macbeth and Lady Macbeth do not use \textit{man} in the same way. Macbeth connects manhood to moral limits: a man may be courageous, but doing more than a man ought to do makes him less than human. Lady Macbeth connects manhood to daring and completion: if Macbeth dares the murder, he proves himself; if he refuses, he is a coward. The argument turns on whether her sense of \textit{man} should replace his.

\textbf{Reading task:} Underline Macbeth's definition of what befits a man. Circle Lady Macbeth's pressure words: \textit{afeard}, \textit{coward}, \textit{dare}, and \textit{man}.

\section*{Passage 3: Act 1, Scene 7 --- Done Does Not Mean Settled}

\textbf{Context:} Macbeth imagines that the murder might be simple if the deed could be finished without consequence. His opening wordplay matters because \textit{done} can mean physically completed, morally settled, or free from future effects.

\begin{playpassage}
\speaker{Macbeth} If it were done when 'tis done, then 'twere well\par
It were done quickly: if the assassination\par
Could trammel up the consequence, and catch\par
With his surcease success; that but this blow\par
Might be the be-all and the end-all here,\par
But here, upon this bank and shoal of time,\par
We'ld jump the life to come. But in these cases\par
We still have judgment here.
\end{playpassage}

\subsection*{Reader Notes}

Macbeth knows that one sense of \textit{done} is too small. The murder can be physically completed in one blow, but that does not mean the consequences are completed, the moral judgment is completed, or the danger is completed. This is not a silly pun. It is a serious logical distinction: an action's immediate completion does not prove that the action is finally settled.

\textbf{Reading task:} Box the repeated word \textit{done}. In the margin, distinguish physical completion from moral or practical settlement.

\section*{How to Use These Passages on the Worksheet}

The worksheet will ask you to apply Week 2 logic to Shakespeare's language. Your job is not to summarize Macbeth's ambition or Lady Macbeth's cruelty. Your job is to define the term, keep senses separate, and test whether an inference survives once the word is clarified.

\textbf{Strong college-level answers} should do three things:

\begin{itemize}[leftmargin=1.5em,itemsep=0.2em]
\item identify the exact word or phrase whose meaning matters;
\item state both senses in ordinary language;
\item explain how the shift affects the reasoning in the scene.
\end{itemize}

\end{document}
